\section{Ubicación del curso e historia del NLP: el Transformer no surgió de la nada}

Esta clase inaugural de la V4 cumple dos tareas. Primero, establece un vocabulario común para las conferencias de los invitados que siguen: attention, large language model, alignment, reasoning y agent. Segundo, presenta una línea de investigación que todavía no se cierra: las capacidades de los modelos se expanden con rapidez, pero la eficiencia en el uso de datos, el aprendizaje continuo, la memoria, la confiabilidad y los límites de seguridad siguen mostrando carencias evidentes. Al leer, conviene separar los ejemplos de productos de la clase de 2024 de los mecanismos más estables: los primeros son evidencia de un momento determinado; solo los segundos son conocimiento transferible.

\lecturefigure{slide-001.jpg}{Portada de CS25 Transformers United V4}{Slide deck oficial de CS25 V4, página 1}

\begin{knowledgebox}{Lectura de la figura: el Overview no es «repasar nombres rápidamente»}
La portada enumera a los cuatro organizadores del curso, y el contenido posterior, en efecto, lo desarrollan distintos ponentes por segmentos. Abarca historia, mecanismos, aplicaciones, debilidades, razonamiento y agents, por lo que funciona más bien como un mapa del curso. Las notas complementan las fórmulas y los límites de los sistemas, pero no reintroducen dashboards, incident runbooks ni procesos de cumplimiento normativo que no aparecieron en la clase.
\end{knowledgebox}

\teachervoice{La apertura explica que el objetivo de CS25 no es repetir los libros de texto, sino que quienes están haciendo investigación expliquen directamente los resultados nuevos, de modo que los estudiantes puedan llevar esas ideas a su propia investigación o establecer colaboraciones. Este enfoque determina que la clase presente lado a lado los consensos, los problemas abiertos y los juicios personales de los ponentes.}

\subsection{Qué espera el curso que se lleven los estudiantes}

Una visión general técnica tiende a degenerar en una lista de nombres de modelos. Los objetivos de aprendizaje que plantea el curso son más concretos: entender cómo funciona el Transformer, entender por qué logró extenderse más allá del NLP, identificar las líneas de investigación de frontera y seguir cuestionando las debilidades que persisten. Este marco exige explicar tanto las condiciones del éxito como los límites del fracaso. Al leer la página de objetivos que sigue, observe primero que los cuatro verbos apuntan, respectivamente, a mechanism, application, research y limitation, en lugar de tomarla como una página promocional del curso.

\lecturefigure{slide-011.jpg}{Los cuatro tipos de objetivos de aprendizaje del curso}{Slide deck oficial de CS25 V4, página 11}

\begin{importantbox}{Los cuatro ejes de esta clase}
\begin{enumerate}
  \item \textbf{Mechanism}: cómo transmiten información attention, self-attention, multi-head y cross-attention;
  \item \textbf{Scale and training}: cómo moldean las capacidades los parámetros, los datos, el alignment y la preference optimization;
  \item \textbf{Adaptation}: cómo la memoria, la recuperación de información, el aprendizaje continuo, el model editing y el reasoning scaffold cubren las carencias del modelo;
  \item \textbf{System}: cómo agent, tool, memory, multi-agent communication y safety convierten el modelo en un sistema ejecutable.
\end{enumerate}
\end{importantbox}

\subsection{Del rule system a la línea de tiempo de attention}

El curso presenta primero una línea de tiempo, no para discutir cuál fue el «primer artículo» de cierta técnica, sino para mostrar qué cuello de botella corregía cada generación de métodos. El NLP rule-based escribe el lenguaje como reglas; el word embedding asigna a cada palabra discreta un vector continuo; las RNN/LSTM procesan secuencias; attention permite que el modelo asigne pesos distintos a distintas posiciones; el Transformer convierte attention en la estructura de cómputo principal.

\lecturefigure{slide-013.jpg}{Línea de tiempo desde los sistemas de reglas hasta el Transformer y los modelos multimodales}{Slide deck oficial de CS25 V4, página 13}

\begin{knowledgebox}{Términos clave: los cinco escalones de la línea de tiempo}
\begin{tabularx}{\textwidth}{L{0.19\textwidth} X}
\toprule
Etapa & Problema principal que resuelve \\
\midrule
Rule-based system & Codifica la sintaxis y la semántica con reglas manuales; es transparente, pero su cobertura es frágil y su costo de transferencia es alto. \\
Word embedding & Expresa la similitud con vectores continuos, lo que permite al modelo compartir la estructura estadística entre palabras. \\
RNN/LSTM & Procesa secuencias de longitud variable con un recurrent state, pero el entrenamiento es secuencial y la información a larga distancia tiende a atenuarse. \\
Attention & Pondera las posiciones de la entrada según su contenido, sin comprimir toda la historia en un único estado. \\
Transformer & Usa attention y redes de propagación hacia adelante como columna vertebral, lo que amplía el margen para el entrenamiento en paralelo y el escalamiento. \\
\bottomrule
\end{tabularx}
\end{knowledgebox}

\lecturefigure{slide-014.jpg}{Ambigüedad, mantenimiento de reglas y generalización: los problemas del NLP temprano}{Slide deck oficial de CS25 V4, página 14}

\begin{warningbox}{La página histórica no respalda que «los métodos antiguos sean completamente inútiles»}
Las reglas, los diccionarios, la recuperación de información, las máquinas de estados finitos y las características tradicionales siguen siendo adecuados para tareas muy restringidas. La ventaja del Transformer está en el aprendizaje de representaciones de extremo a extremo y en el escalamiento, no en invalidar todos los métodos estructurados. Los sistemas de ingeniería suelen usar el learned model junto con reglas, bases de datos y verificadores.
\end{warningbox}

\lecturefigure{slide-015.jpg}{Los métodos de NLP pasan de las reglas escritas a mano a las representaciones aprendidas}{Slide deck oficial de CS25 V4, página 15}

\lecturefigure{slide-016.jpg}{ELIZA muestra las capacidades y los límites de los sistemas de diálogo rule-based}{Slide deck oficial de CS25 V4, página 16}

\begin{knowledgebox}{Lectura de la figura: por qué ELIZA parece «estar conversando»}
ELIZA usa coincidencia de patrones y reescritura con plantillas para insertar fragmentos de la oración del usuario en respuestas predefinidas. Esto muestra que la fluidez superficial no equivale a comprensión interna: si la forma de la interacción es lo bastante adecuada, incluso un conjunto limitado de reglas puede producir una sensación de coherencia. Hoy, al evaluar un LLM, tampoco se puede equiparar la naturalidad del lenguaje con la exactitud factual, la corrección del razonamiento o la confiabilidad de las acciones.
\end{knowledgebox}

\lecturefigure{slide-017.jpg}{Reglas lingüísticas, semantic parsing y las primeras representaciones simbólicas}{Slide deck oficial de CS25 V4, página 17}

\begin{importantbox}{El conocimiento estructural sigue teniendo valor}
El semantic parsing asigna al lenguaje natural una forma lógica o una representación ejecutable. Los LLM modernos pueden reducir las reglas escritas a mano, pero no eliminan la necesidad de estructura: el schema de tool calling, SQL, los programas, los pasos de una demostración y los agent protocols reintroducen, todos ellos, estructuras verificables.
\end{importantbox}

\lecturefigure{slide-018.jpg}{El word embedding asigna a las palabras discretas vectores en un espacio continuo}{Slide deck oficial de CS25 V4, página 18}

\begin{importantbox}{Definición mínima de embedding}
Sea $e_w\in\mathbb{R}^{d}$ el vector de la entrada $w$ del vocabulario. El embedding (incrustación) asigna a cada símbolo discreto un punto de un espacio continuo de $d$ dimensiones, de modo que, bajo cierto objetivo de entrenamiento, las palabras similares obtengan representaciones cercanas. Una medida de similitud habitual es
\[
\operatorname{cos}(e_i,e_j)=\frac{e_i^\top e_j}{\lVert e_i\rVert_2\lVert e_j\rVert_2}.
\]
Aquí $e_i,e_j$ son los vectores de dos tokens u objetos; el numerador es el producto interno y el denominador elimina el efecto de la longitud de los vectores. La cercanía solo indica correlación en el espacio de representación; no equivale automáticamente a una relación factual ni a una relación causal.
\end{importantbox}

\subsection{Resumen de la sección}

La posición histórica del Transformer puede condensarse en una frase: elevó la tarea de «seleccionar la información relevante de una secuencia» a un modo de cómputo principal paralelizable, apilable y escalable. Las capacidades posteriores no provienen de una sola fórmula, sino de la evolución conjunta de architecture, data, optimization, hardware e interface.
